\chapter{Fonctions holomorphes}
\section{Introduction}
Les fonctions holomorphes sont les fonctions complexes dérivables.
\begin{definition}
    
    Soient $\Omega\subseteq \C$ un ouvert et $f:\Omega\to\C$. Soit $z_0\in \Omega$. $f$ est holomorphe en $z_0$ si $\lim_{h\to 0}\frac{f(z_0 + h) - f(z_0)}{h}$ existe. Cette limite est alors notée $f'(z_0)$
    \begin{itemize}
        \item 
        \item $f$ est holomorphe sur $\Omega$ si $f$ est holomorphe en tout point de $\Omega$.
    \end{itemize}
\end{definition}

\begin{exemple}
    Voici quelques exemples de fonctions holomorphes 
    \begin{itemize}
        \item Soit $f(z)=z$. $f$ est holomorphe en tout point car $$\lim_{h\to 0}\frac{f(z+h)- f(z)}{h}= \lim_{h\to 0}\frac{z+h-z}{h}=1\quad \text{ et } \quad f'(z)=1$$ et donc $f'(z) = 1$.
        \item Soit $f(z) = z^2$. $f$ est holomorphe en tout point car $$\lim_{h\to 0}\frac{(z+h)^2-z^2}{h} = \lim_{h\to 0}2z + h = 2z$$
        \item Soit $f(z)= |z|$, $f$ n'est pas holomorphe. car sa restriction à $\R$ n'est pas dérivable.
        \item Soit $f(z)= \overline{z}$. $f$ n'est pas holomorphe. Soit $z\in \C$, on a
        $$
        \lim_{h\to0}\frac{\overline{z+h}- \overline{z}}{h}=\lim_{h\to 0} \frac{\overline{h}}{h}
        $$ n'existe pas car elle vaut $1$ sur la droite réelle et vaut $-1$ sur la droite imaginaire.
    \end{itemize}
\end{exemple}
\begin{proposition}
    Soient $\Omega$ un ouvert de $\C$ et $f,g:\Omega \to \C$. Soit $z_0\in \Omega$.
    \begin{itemize}
        \item Si $f$ et holomorphe en $z_0$ alors $f$ est continue en $z_0$.
        \item Si $f$ et $g$ sont holomorphes en $z_0$ alors $f+g$ est holomorphe en $z_0$ et $(f+g)'(z_0) = f'(z_0) + g'(z_0)$
        \item si $f$ et $g$ sont holomorphes en $z_0$ alors $f g$ est holomorphe en $z_0$ et $(fg)'(z_0)= f'(z_0)g(z_0) + g'(z_0)f(z_0)$
        \item Si $f$ et $g$ sont holomorphes en $z_0$ et qye $g(z_0)\ne 0$ alors $\frac{f}{g}$ est holomorphe en $z_0$ et $\left(\frac{f}{g}\right)'(z_0) = \frac{f'(z_0)g(z_0)-f(z_0)g'(z_0)}{(g(z_0)^2}$
        \item Si $f$ est holomorphe sur $\Omega$ et si $g$ est holomorphe sur une ouvert $V$ tel que $V\subseteq f(\Omega) $ alors $(g\circ f)$ est holomorphe sur $\Omega$ et pour tout $z \in \Omega$ nous avons $(g\circ f)'(z)= g'(f(z))f'(z)$.
    \end{itemize}


\end{proposition}

Cette proposition se prouve de la même manière que les résultats analogues sur $\R$ que nous avons déjà démontré.

\begin{remarque}
    Notons que:
    \begin{itemize}
        \item Tout polynôme complexe est holomorphe sur $\C$.
        \item La fonction logarithme n'est pas holomorphe sur $\C\setminus\{0\}$ puisqu' elle n'est pas continue sur ce même ensemble. Cela aura pour conséquence que $f(z)= \frac{1}{z}$ n'aura pas de primitive sur $\C \setminus \{0\}$.
    \end{itemize}
\end{remarque}
 Tout comme pour les fonctions réelles, nous avons le résultat suivant:

 \begin{proposition}
     Soit $f: \Omega\to \C$ avec $\Omega\subseteq \C$ et soit $z\in\Omega$. La fonction $f$ est holomorphe en $z$ ssi il existe $M_z\in \C$ $$
     f(z+h)= f(z)+M_zh+ R_z(h)
     $$ tel que $\lim_{h\to 0}\frac{R_z(h)}{(h)} = 0$. De plus, on a alors $M_z = f'(z)$.
 \end{proposition}

\begin{proof}
\begin{description}
    \item[$(\alors)$] Supposons $f$ holomorphe en $z$, c'est à dire que 
    $$\lim_{h\to 0}\frac{f(z+h)-f(z)}{h}=f'(z)$$ On a bien, pour $M_z = f'(z)$, que  $$
    \lim_{h\to 0}\frac{R_z(h)}{h}= \lim_{h\to 0}\frac{f(z+h)-f(z)- f(z)h}{h}=\lim_{h\to 0}\frac{f(z+h)-f(z)}{h}- f(z)= 0
    $$
    \item[$(\si)$] Comme $\lim_{h\to 0}\frac{R_z(h)}{h} = \lim_{h\to 0}\frac{f(z+h) - f(z)}{h}-M_z$, on a que $\lim_{h\to 0}\frac{f(z+h) - f(z)}{h}$ existe et vaut $M_z$.
\end{description}
\end{proof}

\newpage

\section{Équations de Cauchy-Riemann}
On sait qu'on peut identifier $\C$ à $\R^2$ en identifiant $z = x + iy$ au couple $(x,y)$. Étant donné donnée un fonction $f:\C\to\C$, on peut l'identifier à une fonction $f_{\R}:\R^2\to\R^2$ en posant
$f_\R(x,y)= (u(x,y), v(x,y))$ où $u(x,y)= \Re(f(x+iy))$ et $v(x+y)= \Im(f(x+iy))$.
\begin{exemple}
Pour illustrer :
\begin{itemize}
    \item Soit $f = \Id_{\C}$. On a que $f_{\R}(x,y) = (x,y)$ d'où $u(x,y) = x$ et $v(x,y) = y$.
    \item Soit $f:\C\to \C; f(z)= f(\overline{z})$ alors $f_{\R}(x,y) = (x,-y)$ et donc $u(x,y) = x$ et $v(x,y) = -y$.
    \item $f(z)= z^2$. alors $f_\R(x,y)=(x^2-y^2, 2xy)$ d'où $u(x,y) = x^2-y^2$ et $v(x,y) = 2xy$..

\end{itemize}
\end{exemple}

\begin{theoreme}[Équations de Cauchy-Riemann]
    Soient $f:\Omega\to \C$ avec $\Omega$ ouvert de $\C$ et $z = x+iy\in \C$.
    $f$ est holomorphe en $z$ ssi 
    \begin{itemize}
        \item  $f_\R$ est différentiable en $(x,y)$
        \item  $\frac{\partial u}{\partial x}(x,y)= \frac{\partial v}{\partial y}(x,y)$ et $\frac{\partial u}{\partial y}(x,y)= \frac{-\partial v}{\partial x}(x,y)$
    \end{itemize}
\end{theoreme}
\begin{proof}
\begin{description}
    \item[$(\alors)$]  Supposons que $f$ est holomorphe en $z$. On a donc $f(z+h) = f(z) + f'(z)  h + R_z(h)$ où $\lim_{h\to 0}\frac{R_z(h)}{h} = 0$
    Autrement dit, 
    \begin{align*}
    f_\R (x+ \Re(h), y + \Im(h))&= f_\R(x,y)+ \left( \Re(f'(z) )\Re(h)- \Im (f'(z))\Im(h), \Re(f'(z)) \Im(h) + \Im(f'(y))\Re(h)  \right) + R_{\R, z}(\Re(h), \Im(h))\\
    &= f_\R(x,y)+  \begin{pmatrix}
        \Re(f'(z)) & - \Im(f'(z)) \\ \Im(f'(z)) & \Re(f'(z))
    \end{pmatrix}\begin{pmatrix}
        \Re(h)\\ \Im(h)
    \end{pmatrix}
+R_{\R, z}(\Re(h), \Im(h)) 
    \end{align*}   
On déduit que $f_\R$ est différentiable en $(x,y)$ et que 
$$
J_{f_\R}(x,y)= \begin{pmatrix}
    \Re(f'(x+i y)) & - \Im(f'(x+iy)) \\ \Im(f'(x+iy)) & \Re(f'(x+iy))
\end{pmatrix}
$$

Or, on sait que 
$$
    J_{f_{\R}} = \begin{pmatrix}
        \frac{\partial u}{\partial x}(x,y) & \frac{\partial u}{\partial y}(x,y)\\
        \frac{\partial v}{\partial x}(x,y) & \frac{\partial v}{\partial x}(x,y)
    \end{pmatrix}
$$
En regardant ces deux matrices, on en déduit que 
$$
\frac{\partial u}{\partial x}(x,y)= \frac{\partial v}{\partial y}(x,y) \quad \text{ et } \quad \frac{\partial u}{\partial y}(x,y)= -\frac{\partial v}{\partial x}(x,y)
$$
\end{description}
   
\end{proof}

Conséquence importante : Si $f$ est holomorphe en $z= x+iy$, on a 
\begin{align*}
    f'(z) &= \frac{\partial u}{\partial x}(x,y)+ i \frac{\partial v}{\partial y}(x,y)\\
    &= \frac{\partial v}{\partial y}(x,y) - i\frac{ \partial u }{\partial y}(x,y) 
\end{align*}
\begin{corollaire}
    Soient $ f: \Omega \to \C $ avec $ \Omega $ ouvert  et $z= x+iy$
    Si $(\frac{\partial u}{\partial x}), \frac{\partial u}{\partial y}, \frac{\partial v}{\partial x}, \frac{\partial v }{\partial y}$ sont continus sur $\Omega$ , alors $f$ est holomorphe en $z$. ie $f$ satisfait :  $$\frac{\partial u}{\partial x}(x,y)= \frac{\partial v}{\partial y}(x,y) \quad \text{ et }\quad \frac{\partial u}{\partial y}(x,y)= \frac{-\partial v}{\partial x}(x,y)$$
\end{corollaire}
\begin{exemple}
Pour illuster :
    \begin{itemize}
        \item $f(z) = z$. On a $u(x,y) = x$ et $v(x,y) = y$. Donc $f$ est holomorphe sur tout le plan complexe.
        \item $f(z) = \overline{z}$
        \item $f(z) = \e^z$. On a $\e^z = \e^{x+\i y} = \e^x(\cos(y)) + \i \e^x\sin(y)$
    \end{itemize}
\end{exemple}
% exo faire avec z^2

\begin{remarque}
Pour rappel :

Soit $f = u + iv : \Omega \to \mathbb{C}$ où $\Omega \subseteq \mathbb{C}$ est ouvert.
Si $\dfrac{\partial u}{\partial x}, \dfrac{\partial u}{\partial y}, \dfrac{\partial v}{\partial x}, \dfrac{\partial v}{\partial y}$
sont continues sur $\{(x,y) : x + iy \in \Omega\}$, alors $f$ est holomorphe en $z = x + iy$ ssi
\[
\left\{
\begin{aligned}
\frac{\partial u}{\partial x}(x,y) &= \frac{\partial v}{\partial y}(x,y) \\[4pt]
\frac{\partial u}{\partial y}(x,y) &= -\frac{\partial v}{\partial x}(x,y)
\end{aligned}
\right.
\qquad \text{(équations de Cauchy--Riemann).}
\]
De plus,
\[
f'(x+iy) = \frac{\partial u}{\partial x}(x,y) + i\,\frac{\partial v}{\partial x}(x,y).
\]
\end{remarque}

\paragraph{Quelques notations supplémentaires.}
Si $f$ est holomorphe en $x+iy$, on peut poser :
\begin{align*}
\frac{\partial f}{\partial x}(x+iy) &:= \frac{\partial u}{\partial x}(x,y) + i\,\frac{\partial v}{\partial x}(x,y) = f'(x+iy), \\[4pt]
\frac{\partial f}{\partial y}(x+iy) &:= \frac{\partial u}{\partial y}(x,y) + i\,\frac{\partial v}{\partial y}(x,y) = i\,f'(x+iy), \\[4pt]
\frac{\partial f}{\partial z}(x+iy) &:= \frac12 \left( \frac{\partial f}{\partial x}(x+iy) - i\,\frac{\partial f}{\partial y}(x+iy) \right) = f'(x+iy), \\[4pt]
\frac{\partial f}{\partial \bar z}(x+iy) &:= \frac12 \left( \frac{\partial f}{\partial x}(x+iy) + i\,\frac{\partial f}{\partial y}(x+iy) \right).
\end{align*}

\begin{proposition}
Soit $f : \Omega \to \mathbb{C}$ où $\Omega$ est un ouvert de $\mathbb{C}$.
Si $\dfrac{\partial u}{\partial x}, \dfrac{\partial u}{\partial y}, \dfrac{\partial v}{\partial x}, \dfrac{\partial v}{\partial y}$
sont continues sur $\{(x,y) : x+iy \in \Omega\}$, alors $f$ est holomorphe en $x+iy \in \Omega$ ssi
\[
\frac{\partial f}{\partial \bar z}(x+iy) = 0 .
\]
\end{proposition}

\begin{proof}
On calcule
\begin{align*}
\frac{\partial f}{\partial \bar z}(x+iy)
&= \frac12 \left( \frac{\partial f}{\partial x}(x+iy) + i\,\frac{\partial f}{\partial y}(x+iy) \right) \\
&= \frac12 \left( \frac{\partial u}{\partial x}(x,y) + i\,\frac{\partial v}{\partial x}(x,y)
   + i \left[ \frac{\partial u}{\partial y}(x,y) + i\,\frac{\partial v}{\partial y}(x,y) \right] \right) \\
&= \frac12 \left( \left[ \frac{\partial u}{\partial x}(x,y) - \frac{\partial v}{\partial y}(x,y) \right]
   + i \left[ \frac{\partial v}{\partial x}(x,y) + \frac{\partial u}{\partial y}(x,y) \right] \right),
\end{align*}
et cette quantité est nulle ssi les deux crochets sont nuls, c'est-à-dire ssi les équations de Cauchy--Riemann sont satisfaites.
\end{proof}

\setcounter{section}{2}
\section{Séries entières et fonctions analytiques}

\subsection{Séries entières}

\begin{definition}[Série entière]
Une \emph{série entière} est une fonction de la forme
\[
f(z) = \sum_{n=0}^{+\infty} a_n (z - z_0)^n ,
\]
où $z_0 \in \mathbb{C}$ et $(a_n)_{n \geq 0}$ est une suite dans $\mathbb{C}$.
\end{definition}

\begin{exemple}
\[
e^z = \sum_{n=0}^{+\infty} \frac{z^n}{n!}, \qquad
\cos(z) = \sum_{n=0}^{+\infty} \frac{(-1)^n}{(2n)!}\, z^{2n}, \qquad
\sin(z) = \sum_{n=0}^{+\infty} \frac{(-1)^n}{(2n+1)!}\, z^{2n+1}.
\]
Dans les trois cas $z_0 = 0$, avec respectivement :
\begin{itemize}
  \item $a_n = \dfrac{1}{n!}$ pour $e^z$ ;
  \item $a_{2n} = \dfrac{(-1)^n}{(2n)!}$, $a_{2n+1} = 0$ pour $\cos$ ;
  \item $a_{2n+1} = \dfrac{(-1)^n}{(2n+1)!}$, $a_{2n} = 0$ pour $\sin$.
\end{itemize}
\end{exemple}

\begin{definition}[Rayon de convergence]
Soit $f(z) = \displaystyle\sum_{n=0}^{+\infty} a_n (z-z_0)^n$. Le \emph{rayon de convergence} de $f$ est donné par
\[
R = \frac{1}{\limsup\limits_{n} \sqrt[n]{|a_n|}} ,
\]
avec les conventions
\[
R = 0 \ \text{ si } \limsup_n \sqrt[n]{|a_n|} = +\infty,
\qquad
R = +\infty \ \text{ si } \limsup_n \sqrt[n]{|a_n|} = 0 .
\]
\end{definition}

\begin{remarque}
Soit $(c_n)_{n \geq 0}$ une suite de réels. On a
\[
\limsup_n c_n = \lim_{n \to +\infty} \Big( \sup_{k \geq n} c_k \Big),
\]
limite qui existe toujours dans $[-\infty, +\infty]$ car la suite $n \mapsto \sup_{k \geq n} c_k$ est décroissante.
Par exemple, si $c_n = (-1)^n$, alors $\limsup_n c_n = 1$.
Si $\lim_n c_n$ existe, alors $\limsup_n c_n = \lim_n c_n$.
\end{remarque}

\begin{exemple}
\begin{itemize}
  \item Pour $f(z) = e^z = \displaystyle\sum_{n=0}^{+\infty} \frac{1}{n!}\, z^n$, on a, par la formule de Stirling
  $n! \sim \sqrt{2\pi n}\,\left(\frac{n}{e}\right)^n$ (c'est-à-dire $\lim_{n\to\infty} \dfrac{n!}{\sqrt{2\pi n}\,(n/e)^n} = 1$) :
  \[
  \limsup_n \sqrt[n]{\frac{1}{n!}}
  = \limsup_n \sqrt[n]{\frac{1}{\sqrt{2\pi n}\,(n/e)^n}}
  = \limsup_n \frac{1}{\sqrt[2n]{2\pi n}\cdot \dfrac{n}{e}}
  = \limsup_n \frac{1}{\underbrace{e^{\frac{1}{2n}\ln(2\pi n)}}_{\to 1}\cdot \dfrac{n}{e}}
  = 0 .
  \]
  Conclusion : $R = +\infty$.
  \item On aura de même $R = +\infty$ pour $\sin$ et $\cos$.
\end{itemize}
\end{exemple}

\begin{proposition}\label{prop:disque}
Soit $f(z) = \displaystyle\sum_{n=0}^{+\infty} a_n (z-z_0)^n$ et soit
\[
D = \Big\{ z \in \mathbb{C} \ \Big|\ \sum_{n=0}^{+\infty} a_n (z-z_0)^n \text{ converge} \Big\}.
\]
Si $R$ est le rayon de convergence, alors
\[
D(z_0, R) \subseteq D \subseteq \overline{D(z_0, R)} .
\]
De plus, sur $D(z_0,R)$ on a convergence absolue, et convergence uniforme sur $\overline{D(z_0, r)}$ pour tout $r < R$.
\end{proposition}

\begin{proof}
Soit $z \in D(z_0, R)$. Montrons que $\sum_{n=0}^{+\infty} a_n (z-z_0)^n$ converge absolument, c'est-à-dire que
$\sum_{n=0}^{+\infty} |a_n|\,|z-z_0|^n$ converge.
Par le critère de Cauchy, si $\limsup_n \sqrt[n]{|a_n|\,|z-z_0|^n} < 1$ alors on aura convergence de la série. Or
\[
\limsup_n \sqrt[n]{|a_n|\,|z-z_0|^n} < 1
\iff \Big( \limsup_n \sqrt[n]{|a_n|} \Big)\,|z-z_0| < 1
\iff |z-z_0| < \frac{1}{\limsup_n \sqrt[n]{|a_n|}} = R
\]
(si $\limsup_n \sqrt[n]{|a_n|} \neq 0, +\infty$ ; les cas limites se traitent de même).

De même, si $z \notin \overline{D(z_0,R)}$ alors $|z-z_0| > R$, et par le critère de Cauchy on peut conclure que
$\sum_{n=0}^{+\infty} a_n (z-z_0)^n$ diverge.

Pour finir, si on considère $r < R$, on aura convergence uniforme sur $\overline{D(z_0,r)}$ car
\[
\sum_{n=0}^{+\infty} \sup_{z \in \overline{D(z_0,r)}} \big( |a_n|\,|z-z_0|^n \big) = \sum_{n=0}^{+\infty} |a_n|\, r^n ,
\]
qui converge puisque $\limsup_n \sqrt[n]{|a_n|\, r^n} = r \cdot \limsup_n \sqrt[n]{|a_n|} = \dfrac{r}{R} < 1$.
\end{proof}

On peut en fait montrer que toute série entière $f(z) = \sum_{n=0}^{+\infty} a_n (z-z_0)^n$ est holomorphe sur $D(z_0,R)$.

\begin{theoreme}\label{thm:derivation}
Soit $f(z) = \displaystyle\sum_{n=0}^{+\infty} a_n (z-z_0)^n$ et soit $R$ le rayon de convergence.
Alors $f$ est holomorphe sur $D(z_0,R)$ et, pour tout $z \in D(z_0,R)$,
\[
f'(z) = \sum_{n=1}^{+\infty} n\, a_n (z-z_0)^{n-1} = \sum_{n=0}^{+\infty} (n+1)\, a_{n+1} (z-z_0)^n .
\]
De plus, le rayon de convergence de $f'$ vaut $R$.
\end{theoreme}

\begin{proof}
\textbf{Étape 1 : rayon de convergence.}
Montrons que $\sum_{n=0}^{+\infty} (n+1)\, a_{n+1} (z-z_0)^n$ a un rayon de convergence égal à $R$.
Remarquons que cette série converge ssi $\sum_{n=0}^{+\infty} (n+1)\, a_{n+1} (z-z_0)^{n+1}$ converge. Or
\[
\limsup_n \sqrt[n+1]{(n+1)\,|a_{n+1}|}
= \limsup_n \underbrace{\sqrt[n+1]{n+1}}_{\to 1} \cdot \sqrt[n+1]{|a_{n+1}|}
= \limsup_n \sqrt[n]{|a_n|} = \frac{1}{R} .
\]

\textbf{Étape 2 : dérivabilité.}
Il reste à montrer que pour $z \in D(z_0,R)$,
\[
(\ast)\qquad \lim_{h \to 0} \left| \frac{f(z+h) - f(z)}{h} - \sum_{n=0}^{+\infty} (n+1)\, a_{n+1} (z-z_0)^n \right| = 0 .
\]
Soit $z \in D(z_0,R)$. Quitte à travailler avec des $h$ très petits, on peut considérer $r < R$ tel que
$z$ et $z+h$ appartiennent à $D(z_0,r)$. On a
\[
(\ast) = \lim_{h \to 0} \left|
\frac{\displaystyle\sum_{n=0}^{+\infty} a_n (z+h-z_0)^n - \sum_{n=0}^{+\infty} a_n (z-z_0)^n}{h}
- \sum_{n=0}^{+\infty} (n+1)\, a_{n+1} (z-z_0)^n \right| ,
\]
et, pour n'importe quel choix de $N$,
\begin{multline*}
(\ast) = \lim_{h \to 0} \Bigg|
\frac{\displaystyle\sum_{n=0}^{N} a_n (z+h-z_0)^n - \sum_{n=0}^{N} a_n (z-z_0)^n}{h}
- \sum_{n=0}^{N-1} (n+1)\, a_{n+1} (z-z_0)^n \\
+ \underbrace{\frac{\displaystyle\sum_{n=N+1}^{+\infty} a_n (z+h-z_0)^n - \sum_{n=N+1}^{+\infty} a_n (z-z_0)^n}{h}}_{(1)}
- \underbrace{\sum_{n=N}^{+\infty} (n+1)\, a_{n+1} (z-z_0)^n}_{(2)}
\Bigg| .
\end{multline*}
Les deux premiers termes forment une somme finie dont la différence tend vers $0$ quand $h \to 0$
(dérivation terme à terme d'un polynôme). Remarquons de plus que :
\begin{itemize}
  \item $(2)$ est arbitrairement petit si $N$ est choisi suffisamment grand (reste d'une série convergente, par l'étape 1) ;
  \item $(1) = \displaystyle\sum_{n=N+1}^{+\infty} a_n \, \frac{(z+h-z_0)^n - (z-z_0)^n}{h}$.
  Cependant, en posant $b = z+h-z_0$ et $c = z-z_0$,
  \[
  b^n - c^n = (b-c)\big( b^{n-1} + c\,b^{n-2} + c^2 b^{n-3} + \dots + c^{n-2} b + c^{n-1} \big),
  \]
  et donc
  \[
  (1) = \sum_{n=N+1}^{+\infty} a_n \cdot \frac{h \cdot \overbrace{\big( (z+h-z_0)^{n-1} + \dots + (z-z_0)^{n-1} \big)}^{|\cdot| \,\leq\, n\, r^{n-1}}}{h} .
  \]
  Or $\displaystyle\sum_{n=0}^{+\infty} |a_n|\, n\, r^{n-1}$ converge (car $r < R$ et le rayon de la série dérivée vaut $R$),
  et donc $(1)$ est arbitrairement petit si $N$ est choisi suffisamment grand, uniformément en $h$.
\end{itemize}
Ceci conclut la preuve de $(\ast)$.
\end{proof}

\begin{corollaire}
Soit $f(z) = \displaystyle\sum_{n=0}^{+\infty} a_n (z-z_0)^n$ avec $R$ son rayon de convergence.
Alors $f$ est holomorphe à tout ordre sur $D(z_0,R)$ et, pour tout $k \geq 1$,
\[
f^{(k)}(z) = \sum_{n=k}^{+\infty} n(n-1)\cdots(n-k+1)\, a_n (z-z_0)^{n-k}
= \sum_{n=0}^{+\infty} (n+k)\cdots(n+1)\, a_{n+k} (z-z_0)^n .
\]
De plus, on a
\[
a_k = \frac{f^{(k)}(z_0)}{k!} .
\]
\end{corollaire}

\begin{proof}
Le fait que $f$ est holomorphe à tout ordre sur $D(z_0,R)$ découle du Théorème~\ref{thm:derivation} :
la dérivée d'une série entière de rayon $R$ redonne une série entière de rayon $R$, et on itère.
De plus, comme
\[
f^{(k)}(z) = \sum_{n=0}^{+\infty} (n+k)\cdots(n+1)\, a_{n+k} (z-z_0)^n ,
\]
en évaluant en $z = z_0$ seul le terme $n = 0$ subsiste :
\[
f^{(k)}(z_0) = \underbrace{k \cdot (k-1) \cdots 1}_{k!}\ a_k .
\qedhere
\]
\end{proof}